% GENERATED FILE. Edit canonical lessons, then run npm run generate:books.
% canonical-source-hash: fnv1a64:31e0cfc1d041e7db
% canonical-lessons: ZH-C161-dami, ZH-W161-mi-rice, ZH-C161-fen, ZH-C161-chaofan, ZH-C161-niurou, ZH-C161-jirou

\chapter{Rice and Meat}
\label{ch:zh-rice-and-meat}
\hlchaptermodality{161}

\begin{chapteropening}
\textbf{By the end of this chapter:} \emph{I can ask for rice, a portion, fried rice, beef and chicken.}

Complete the last lesson of chapter 161: I can ask for rice, a portion, fried rice, beef and chicken.
\end{chapteropening}

\section[\texorpdfstring{\zh{大米} (dàmǐ)}{大米 (dàmǐ)}]{\zh{大米} (dàmǐ) — uncooked rice; rice grain}

\label{lesson:ZH-C161-dami}



\begin{quote}
Before the new one: say the Mandarin for to get the jobs done, then the Mandarin for to look after.
\end{quote}

\subsection*{You'll want to know: \zh{大米}}
\textbf{\zh{大米}} — \emph{dàmǐ} — \textquotedblleft{}uncooked rice; rice grain\textquotedblright{}.

\begin{quote}
\textbf{\zh{我去买大米}} \emph{wǒ qù mǎi dàmǐ} — I'm going to buy rice.
\end{quote}

The new character, in pieces: \zh{米} — six strokes: two short strokes slanting in toward the middle, a horizontal, a vertical through it, then from the crossing a stroke falling to the left and one falling to the right. The next lesson puts it in your hand.

\begin{cousinweb}
\textbf{\zh{大米}} (\emph{dàmǐ}) is built from \zh{大} 'big' and \zh{米} 'grain', both inherited from Middle Chinese, \emph{daH} and \emph{mejX} (Guangyun, in Baxter's transcription), which glosses the second as the seed of a cereal.
\end{cousinweb}

\begin{grammarlens}[title={What you've built}]
One more word for asking for what I want.
\end{grammarlens}

\subsection*{Your turn}
\begin{itemize}
\raggedright
  \item \emph{Say it:} \emph{dàmǐ}
  \item \emph{Say it:} \emph{dàmǐ}, once more
  \item \emph{Say it:} \emph{wǒ qù mǎi dàmǐ} — I'm going to buy rice.
  \item \emph{Recall:} say the Mandarin for to get the jobs done, then the Mandarin for to look after, then say \emph{dàmǐ} again
\end{itemize}

\begin{tcolorbox}[breakable,colback=teal!4,colframe=teal!35!black,title={Before you move on}]
What does \zh{大米} mean? (\textquotedblleft{}Uncooked rice; rice grain\textquotedblright{}.) Say it once more.
\end{tcolorbox}
\section[\texorpdfstring{\zh{米} (mǐ)}{米 (mǐ)}]{\zh{米} — a cross with two strokes above and two below}

\label{lesson:ZH-W161-mi-rice}



\begin{quote}
The word first, then the character inside it.

\begin{itemize}
\raggedright
  \item \emph{Recall:} say \emph{dàmǐ} — \textbf{R1}, one lesson back, before it decays
\end{itemize}

\begin{quote}
\zh{来} \zh{大}
\end{quote}

\emph{Read it:} each one aloud, with nothing said first — \emph{lái}, \emph{dà}
\end{quote}

\subsection*{Script — a cross with two strokes above and two below}
\hlblockfigure{\detokenize{figures/ZH-W161-mi-rice-filmstrip.pdf}}{How \zh{米} is written, stroke by stroke}

\begin{quote}
\zh{米}
\end{quote}

\textbf{One piece, six strokes, built round a cross.}

\textbf{Two short strokes first,} slanting in toward the middle: a dot from the upper left, then a short stroke from the upper right falling to the left.

\textbf{The cross.} A long horizontal, then a vertical down through its middle.

\textbf{The legs, from the crossing.} A stroke falling to the lower left, then one falling to the lower right.

2 + 2 + 2 = six.

\subsection*{Your turn}
\begin{itemize}
\raggedright
  \item \emph{Write it:} \zh{米}, all six, with the model in front of you
  \item \emph{Cover:} write it again from memory, and count the pieces, not the strokes
\end{itemize}

\begin{tcolorbox}[breakable,colback=teal!4,colframe=teal!35!black,title={Before you move on}]
How many strokes in \zh{米}? (\textbf{Six}.) Which two come first? (\textbf{The two short strokes above the cross}.) Which come last? (\textbf{The two legs, left, then right}.)

Source: \href{https://raw.githubusercontent.com/chanind/hanzi-writer-data/68d10a4b21150cae5e1ebbd223eed289cf32d90c/data/\%E7\%B1\%B3.json}{Hanzi Writer Data \zh{米}, pinned stroke paths and medians}.
\end{tcolorbox}
\section[fèn]{fèn — measure word for a serving; portion}

\label{lesson:ZH-C161-fen}



\begin{quote}
Before the new one: say the Mandarin for to look after, then the Mandarin for uncooked rice.
\end{quote}

\subsection*{You'll want to know: fèn}
\emph{fèn} — \textquotedblleft{}measure word for a serving; portion\textquotedblright{}.

\begin{quote}
\emph{wǒ yào liǎng fèn mǐfàn} — Two portions of rice, please.
\end{quote}

Its characters come later; for now it is a word for the ear.

\begin{cousinweb}
\emph{fèn} has no older form of this word in any dictionary we use: the Guangyun lists the character only as \emph{pin}, glossed 'refined', which is neither this sound nor this sense. Its history is not settled.
\end{cousinweb}

\begin{grammarlens}[title={What you've built}]
One more word for asking for what I want.
\end{grammarlens}

\subsection*{Your turn}
\begin{itemize}
\raggedright
  \item \emph{Say it:} \emph{fèn}
  \item \emph{Say it:} \emph{fèn}, once more
  \item \emph{Say it:} \emph{wǒ yào liǎng fèn mǐfàn} — Two portions of rice, please.
  \item \emph{Recall:} say the Mandarin for to look after, then the Mandarin for uncooked rice, then say \emph{fèn} again
\end{itemize}

\begin{tcolorbox}[breakable,colback=teal!4,colframe=teal!35!black,title={Before you move on}]
What does fèn mean? (\textquotedblleft{}Measure word for a serving; portion\textquotedblright{}.) Say it once more.
\end{tcolorbox}
\section[chǎofàn]{chǎofàn — fried rice}

\label{lesson:ZH-C161-chaofan}



\begin{quote}
Before the new one: say the Mandarin for uncooked rice, then the Mandarin for measure word for a serving.
\end{quote}

\subsection*{You'll want to know: chǎofàn}
\emph{chǎofàn} — \textquotedblleft{}fried rice\textquotedblright{}.

\begin{quote}
\emph{tā yào le yí fèn chǎofàn} — He ordered a portion of fried rice.
\end{quote}

Its characters come later; for now it is a word for the ear.

\begin{cousinweb}
\emph{chǎofàn} is built from \emph{chǎo} 'to stir-fry' and \zh{饭} 'cooked rice', both inherited from Middle Chinese, \emph{tsrhaewX} and \emph{bjonH} (Guangyun, in Baxter's transcription), which glosses the first 'to roast dry'.
\end{cousinweb}

\begin{grammarlens}[title={What you've built}]
One more word for asking for what I want.
\end{grammarlens}

\subsection*{Your turn}
\begin{itemize}
\raggedright
  \item \emph{Say it:} \emph{chǎofàn}
  \item \emph{Say it:} \emph{chǎofàn}, once more
  \item \emph{Say it:} \emph{tā yào le yí fèn chǎofàn} — He ordered a portion of fried rice.
  \item \emph{Recall:} say the Mandarin for uncooked rice, then the Mandarin for measure word for a serving, then say \emph{chǎofàn} again
\end{itemize}

\begin{tcolorbox}[breakable,colback=teal!4,colframe=teal!35!black,title={Before you move on}]
What does chǎofàn mean? (\textquotedblleft{}Fried rice\textquotedblright{}.) Say it once more.
\end{tcolorbox}
\section[niúròu]{niúròu — beef}

\label{lesson:ZH-C161-niurou}



\begin{quote}
Before the new one: say the Mandarin for measure word for a serving, then the Mandarin for fried rice.
\end{quote}

\subsection*{You'll want to know: niúròu}
\emph{niúròu} — \textquotedblleft{}beef\textquotedblright{}.

\begin{quote}
\emph{zhège niúròu hěn hǎochī} — This beef is delicious.
\end{quote}

Its characters come later; for now it is a word for the ear.

\begin{cousinweb}
\emph{niúròu} is built from \emph{niú} 'ox' and \emph{ròu} 'meat', both inherited from Middle Chinese, \emph{ngjuw} and \emph{nyuwk} (Guangyun, in Baxter's transcription).
\end{cousinweb}

\begin{grammarlens}[title={What you've built}]
One more word for asking for what I want.
\end{grammarlens}

\subsection*{Your turn}
\begin{itemize}
\raggedright
  \item \emph{Say it:} \emph{niúròu}
  \item \emph{Say it:} \emph{niúròu}, once more
  \item \emph{Say it:} \emph{zhège niúròu hěn hǎochī} — This beef is delicious.
  \item \emph{Recall:} say the Mandarin for measure word for a serving, then the Mandarin for fried rice, then say \emph{niúròu} again
\end{itemize}

\begin{tcolorbox}[breakable,colback=teal!4,colframe=teal!35!black,title={Before you move on}]
What does niúròu mean? (\textquotedblleft{}Beef\textquotedblright{}.) Say it once more.
\end{tcolorbox}
\section[jīròu]{jīròu — chicken meat; chicken (as food)}

\label{lesson:ZH-C161-jirou}



\begin{quote}
Before the new one: say the Mandarin for fried rice, then the Mandarin for beef.
\end{quote}

\subsection*{You'll want to know: jīròu}
\emph{jīròu} — \textquotedblleft{}chicken meat; chicken (as food)\textquotedblright{}.

\begin{quote}
\emph{nǐ yào niúròu háishi jīròu} — Would you like beef or chicken?
\end{quote}

Its characters come later; for now it is a word for the ear.

\begin{cousinweb}
\emph{jīròu} is built from \emph{jī} 'chicken' and \emph{ròu} 'meat', both inherited from Middle Chinese, \emph{kej} and \emph{nyuwk} (Guangyun, in Baxter's transcription); the dictionary quotes an older book calling the first the creature that knows the hour.
\end{cousinweb}

\begin{grammarlens}[title={What you've built}]
One more word for asking for what I want.
\end{grammarlens}

\subsection*{Your turn}
\begin{itemize}
\raggedright
  \item \emph{Say it:} \emph{jīròu}
  \item \emph{Say it:} \emph{jīròu}, once more
  \item \emph{Say it:} \emph{nǐ yào niúròu háishi jīròu} — Would you like beef or chicken?
  \item \emph{Recall:} say the Mandarin for fried rice, then the Mandarin for beef, then say \emph{jīròu} again
\end{itemize}

\begin{tcolorbox}[breakable,colback=teal!4,colframe=teal!35!black,title={Before you move on}]
What does jīròu mean? (\textquotedblleft{}Chicken meat; chicken (as food)\textquotedblright{}.) Say it once more.
\end{tcolorbox}
